\documentclass[10pt,a4paper]{article}
\usepackage[margin=1in]{geometry}
\usepackage{hyperref}
\hypersetup{
    colorlinks=true,
    linkcolor=blue,
    urlcolor=blue,
    citecolor=blue
}
\usepackage{enumitem}
\usepackage{titlesec}
\usepackage{parskip}
\usepackage{fontawesome5}  
\usepackage[spanish]{babel}

\titleformat{\section}{\large\bfseries}{}{0em}{}[\titlerule]

\begin{document}

\begin{center}
    {\Huge \textbf{Fernando Gutiérrez Canales}}\\
    {\Large Científico de Datos | Doctorado (PhD)}\\    
    León, México \quad | \quad \href{https://www.linkedin.com/in/fernando-guti%C3%A9rrez-canales-804876343/}{LinkedIn} \quad | \quad +52 1 477 55 18 832 \quad | \quad \href{mailto:carl.cfgc@gmail.com}{carl.cfgc@gmail.com} \\
     \href{https://github.com/Fernando-Canales}{GitHub: github.com/Fernando-Canales}
\end{center}

\vspace{1em}

\section*{Resumen Profesional}
Científico de Datos con Doctorado en Astrofísica y más de 5 años de experiencia diseñando \textit{pipelines} de datos escalables, evaluando Modelos de Lenguaje Grande (LLMs) y aplicando análisis estadístico avanzado a conjuntos de datos masivos y complejos. Experto en Python, C y \textit{frameworks} de \textit{machine learning}. Capacidad comprobada para traducir investigaciones científicas rigurosas en soluciones algorítmicas optimizadas, cerrando la brecha entre la arquitectura de datos altamente técnica y los resultados de negocio de gran impacto.

\section*{Experiencia Profesional}
\textbf{Mercor} \hfill \textbf{Göttingen (Remoto)} \\
\textit{Ingeniero de Prompts de IA (Entrenamiento y Evaluación de LLMs)} \hfill \textit{Oct 2025 -- Feb 2026}
\begin{itemize}[leftmargin=1.5em]
	\item Desarrollé conjuntos de datos lógicos y analíticos complejos (prompts adversarios) para entrenar y refinar (\textit{fine-tune}) Modelos de Lenguaje Grande (LLMs), probando y mejorando efectivamente su capacidad de razonamiento en un 30\%.
\end{itemize}

\textbf{Observatorio de París \& Instituto Max Planck de Investigaciones del Sistema Solar} \hfill \textbf{París y Göttingen} \\
\textit{Científico de Datos e Investigador Doctoral} \hfill \textit{Mar 2022 -- Jul 2025}
\begin{itemize}[leftmargin=1.5em]
    \item Diseñé e implementé \textit{pipelines} de datos automatizados en Python para procesar conjuntos de datos astronómicos masivos, utilizando métodos estadísticos avanzados y análisis de series de tiempo para estimar eficiencias de detección de exoplanetas.
    \item Integré módulos de C y Bash para optimizar el rendimiento algorítmico, reduciendo drásticamente la carga computacional y el tiempo de manejo manual de datos.
    \item Colaboré con científicos de más de 15 instituciones internacionales para estandarizar arquitecturas de datos complejas y garantizar la estricta integridad de los datos para la misión espacial PLATO.
\end{itemize}

\textbf{ESTEC, Agencia Espacial Europea (ESA)} \hfill \textbf{Noordwijk, Países Bajos}\\
\textit{Científico de Datos (Prácticas)}  \hfill \textit{Jun 2023 -- Ago 2023}
\begin{itemize}[leftmargin=1.5em]
    \item Desarrollé \textit{scripts} automatizados utilizando Python para extraer, validar y analizar estadísticamente parámetros complejos de detectores fotométricos, mejorando la confiabilidad de los \textit{pipelines} de datos críticos para la misión.
\end{itemize}

\section*{Proyectos Destacados en Ciencia de Datos}
\textbf{Modelado Predictivo y Análisis Estadístico (TripleTen)} \hfill \textbf{Python, Scikit-learn, Pandas} \\
\textit{Científico de Datos} \hfill \textit{Marzo 2026 -- Abril 2026}
\begin{itemize}[leftmargin=1.5em]
    \item Desarrollé un modelo de \textit{machine learning} utilizando Scikit-learn para predecir la tasa de abandono (\textit{churn}) de clientes, logrando una precisión del 85\%.
    \item Llevé a cabo rigurosas pruebas A/B y pruebas de hipótesis estadísticas para validar el impacto del modelo en las estrategias de retención del negocio.
\end{itemize}

\textbf{Arquitectura de Datos para Clínica Oftalmológica} \hfill \textbf{León, México} \\
\textit{Consultor de Datos Principal (Proyecto Independiente)} \hfill \textit{2021 -- 2022}
\begin{itemize}[leftmargin=1.5em]
    \item Ejecuté el proceso completo de limpieza y estructuración de datos (\textit{data wrangling}) sobre años de registros físicos y digitales no estructurados, construyendo una base de datos relacional centralizada y robusta.
    \item Diseñé \textit{scripts} automatizados para extraer \textit{insights} accionables, reduciendo el tiempo de recuperación de información y optimizando la asignación de recursos de la clínica.
\end{itemize}

\section*{Habilidades Técnicas}
\textbf{Programación y Stack Principal:} Python (Pandas, NumPy, Scikit-learn, SciPy), C, Bash, SQL, R. \\
\textbf{Ciencia de Datos e IA:} Machine Learning, Ingeniería de Prompts para LLMs, Modelado Estadístico, Análisis de Series de Tiempo, Analítica Predictiva, Optimización de Algoritmos. \\
\textbf{Ingeniería de Datos y Herramientas:} Pipelines Automatizados, Arquitectura de Datos, Git, Docker, Linux, Tableau.

\section*{Educación}
\textbf{Certificación de Bootcamp en Análisis de Datos} \hfill \textit{TripleTen} \hfill \textit{2026}\\
\textbf{Doctorado en Astrofísica} \hfill \textit{Observatorio de París \& MPS Göttingen} \hfill \textit{2025}\\
\textbf{Maestría en Astrofísica} \hfill \textit{Universidad de Guanajuato, México} \hfill \textit{2021}\\
\textbf{Licenciatura en Física} \hfill \textit{Universidad de Guanajuato, México} \hfill \textit{2019}

\section*{Logros y Publicaciones}
\begin{itemize}[leftmargin=1.5em]
    \item Coautor de \href{https://ui.adsabs.harvard.edu/search/fq...}{publicaciones revisadas por pares} aplicando métodos estadísticos avanzados para el análisis de datos fotométricos.
    \item Seleccionado para presentar hallazgos de investigación predictiva en conferencias internacionales de primer nivel (EAS Padua 2024, PLATO Week Praga 2023).
    \item Beneficiario de becas académicas completas y graduado con Mención Honorífica en los tres grados universitarios (2019, 2021, 2025).
\end{itemize}

\end{document}
